\begin{afterword}
\summaryhead

The post defines a cognitive bias. It is an obstacle to truth that comes not from the cost of information or from limited computing power, but from the shape of our own mental machinery. That shape might have evolved to win arguments, or to believe what others believe, or it might be the side effect of a useful shortcut such as the availability heuristic.

Biases are set apart from mistakes, which are errors of learned content and are easier to fix; from the effects of brain damage; and from absorbed culture. Biases come from machinery all humans share. Plato's ignorance of relativity was not a bias, but a belief in philosopher-kings would be one if it came from a universal instinct for self-promotion.

The author then says the term does not matter much, and can be dropped if unhelpful. We fight biases not out of duty but because they stand between us and the truth, which we need and are curious about.

\responsehead

The post is a definition, and the definition is harder to use than it looks.

It sorts errors by their cause. An error is a bias if it comes from the shape of universal machinery, and not a bias if it comes from missing information, limited computing power, upbringing, culture or damage. The Plato example shows what this means, but it also shows the difficulty: whether Plato's belief was a bias depends on whether it came from an instinct, from his father, or, in the post's joke, from sniffing glue as a child, and the post gives no way to find out. The same goes for ``humanly universal'': the post does not say how one would establish it.

One line inside the definition is not drawn at all. Errors from ``limited computing power'' are excluded, yet the post counts the errors of a shortcut like availability as biases, and a shortcut is used because the full calculation costs too much. The post does not say how to tell the two causes apart.

Two claims come without support. ``In cognitive science'' presents the distinction between biases and mistakes as the field's, with no source, and mistakes are said to be ``much easier to correct'' with no evidence.

After five paragraphs on the definition, the author says the project does not depend on it, and the closing paragraphs explain why biases matter rather than what they are.

In short: the post defines a bias by its cause, gives no way to find the cause of a particular error, and then says the definition does not matter much.
\end{afterword}
